\section*{章末练习}

\begin{longtable}{L{0.08\linewidth}L{0.32\linewidth}L{0.5\linewidth}}
\toprule
题号 & 类型 & 题目 \\
\midrule
1 & 电平判断 & 某输入规定 $V_{IL,max}=0.8$ V、$V_{IH,min}=2.0$ V。分别判断 0.5 V、1.4 V 和 2.7 V 属于哪一区间，并说明中间区为何不能当作第三种逻辑值。 \\
2 & 噪声余量 & 驱动器保证 $V_{OH,min}=2.7$ V、$V_{OL,max}=0.4$ V，接收器指标同题 1。计算高、低电平噪声余量。 \\
3 & 器件路径 & 用受控导通路径解释 CMOS 反相器输入为 0 和 1 时输出为何分别趋近电源与地。 \\
4 & 功耗 & 一个节点每秒翻转次数加倍、负载电容不变，动态功耗如何变化？若电压降为原来的 0.8 倍呢？ \\
5 & 边沿 & 为什么慢边沿会增加门电路短路电流和对噪声的敏感时间？施密特触发输入能改善什么？ \\
6 & 排错 & 板上数字输入偶发误判，按电源、阈值、边沿、串扰和测量参考地列出检查顺序。 \\
\bottomrule
\end{longtable}

\section*{参考解答与要点}

\begin{longtable}{L{0.08\linewidth}L{0.32\linewidth}L{0.5\linewidth}}
\toprule
题号 & 类型 & 解答要点 \\
\midrule
1 & 电平判断 & 0.5 V 保证为低，2.7 V 保证为高，1.4 V 位于未规定区；该区的输出受器件、噪声和时间影响，不能作为稳定逻辑状态。 \\
2 & 噪声余量 & 高电平为 $2.7-2.0=0.7$ V，低电平为 $0.8-0.4=0.4$ V。 \\
3 & 器件路径 & 输入低时上拉网络导通、下拉截止，输出电容向电源充电；输入高时相反，输出经下拉网络放电到地。 \\
4 & 功耗 & 动态功耗近似正比于翻转率，次数加倍则加倍；电压为 0.8 倍时电压项使功耗乘 $0.8^2=0.64$。 \\
5 & 边沿 & 慢边沿让上下管同时部分导通更久，也延长输入停留在阈值附近的时间；施密特触发以滞回提高慢变或带噪输入的稳定性。 \\
6 & 排错 & 先确认电源和参考地，再量接收端实际高低电平与余量，检查上升下降时间、反射/串扰和探头接地，最后核对输入类型与阈值配置。 \\
\bottomrule
\end{longtable}
